\section{USAMO 2012}
        \begin{enumerate}
        	\item (\textit{USAMO 2012}) Find all integers $n \ge 3$ such that among any $n$ positive real numbers $a_1$, $a_2$, $\dots$, $a_n$ with

 \[\max(a_1, a_2, \dots, a_n) \le n \cdot \min(a_1, a_2, \dots, a_n),\]
there exist three that are the side lengths of an acute triangle.


 



\textbf{Solution}

Without loss of generality, assume that the set $\{a\}$ is ordered from least to greatest so that the bounding condition becomes $a_n \le n \cdot a_1.$ Now set $b_i \equiv \frac{a_i}{a_1},$ and since a triangle with sidelengths from $\{a\}$ will be similar to the corresponding triangle from $\{b\},$ we simply have to show the existence of acute triangles in $\{b\}.$ Note that $b_1 = 1$ and for all $i$, $b_i \le n.$

 Now three arbitrary sidelengths $x$, $y$, and $z$, with $x \le y \le z,$ will form a valid triangle if and only if $x+y>z.$ Furthermore, this triangle will be acute if and only if $x^2 + y^2 > z^2.$ However, the first inequality can actually be inferred from the second, since $x+y>z \longleftrightarrow x^2 + y^2 +2xy > z^2$ and $2xy$ is trivially greater than $0.$ So we just need to find all $n$ such that there is necessarily a triplet of $b$'s for which $b_i^2 + b_j^2 > b_k^2$ (where $b_i < b_j < b_k$).

 We now make another substitution: $c_i \equiv b_i ^2.$ So $c_1 = 1$ and for all $i$, $c_i \le n^2.$ Now we examine the smallest possible sets $\{c\}$ for small $n$ for which the conditions of the problem are not met. Note that by smallest, we mean the set whose greatest element is as small as possible. If $n=3$, then the smallest possible set, call it $\{s_3\},$ is trivially $\{1,1,2\}$, since $c_1$ and $c_2$ are obviously minimized and $c_3$ follows as minimal. Using this as the base case, we see inductively that in general $\{s_n\}$ is the set of the first $n$ Fibonacci numbers. To show this note that if $\{s_n\} = \{F_0, F_1, ... F_n\}$, then $\{s_{n+1}\} = \{F_0, F_1, ... F_n, c_{n+1}\}.$ The smallest possible value for $c_{n+1}$ is the sum of the two greatest values of $\{s_n\}$ which are $F_{n-1}$ and $F_n$. But these sum to $F_{n+1}$ so $\{s_{n+1}\} = \{F_0, F_1, ... F_{n+1}\}$ and our induction is complete.

 Now since we know that the Fibonacci set is the smallest possible set which does not satisfy the conditions of the problem, then any set $\{c\}$ whose greatest term is less than $F_{n-1}$ must satisfy the conditions. And since $\{c\}$ is bounded between $1$ and $n^2$, then the conditions of the problem are met if and only if $F_{n-1} > n^2$. The first $n$ for which this restriction is satisfied is $n=13$ and the exponential behavior of the Fibonacci numbers ensure that every $n$ greater than $13$ will also satisfy this restriction. So the final solution set is $\boxed{\{n \ge 13\}}$.

 



\textbf{Solution 2}

\textbf{Outline:}

 1. Define the Fibonacci numbers to be $F_1 = F_2 = 1$ and $F_k = F_{k-1} + F_{k-2}$ for $k \ge 3$.

 2. If the chosen $n$ is such that $F_n \le n^2$, then choose the sequence $a_n$ such that $a_k = \sqrt{F_k}$ for $1 \le k \le n$. It is easy to verify that such a sequence satisfies the condition that the largest term is less than or equal to $n$ times the smallest term. Also, because for any three terms $x = \sqrt{F_a}, y = \sqrt{F_b}, z = \sqrt{F_c}$ with $a<b<c$, $x^2 + y^2 = F_a + F_b \le F_{b-1} + F_b = F_{b+1} \le F_c = z^2$, x, y, z do not form an acute triangle. Thus, all $n$ such that $F_n \le n^2$ do not work.

 3. It is easy to observe via a contradiction argument that all $n$ such that $F_n > n^2$ produce an acute triangle. (If, without loss of generality, $a_n$ is an increasing sequence, such that no three (in particular, consecutive) terms form an acute triangle, then $a_2^2 \ge F_1a_1^2, a_3^2 \ge a_2^2 + a_1^2 \ge F_2a^2$, and by induction $a_n^2 > F_na_1^2$, a contradiction to the condition's inequality.)

 4. Note that $F_{12} = 144 = 12^2$ and $F_{13} = 233 > 169 = 13^2$. It is easily to verify through strong induction that all $n$ greater than 12 make $F_n > n^2$. Thus, $\boxed{n \ge 13}$ is the desired solution set.

 


 


	\item (\textit{USAMO 2012}) A circle is divided into $432$ congruent arcs by $432$ points.  The points are colored in four colors such that $108$ points are colored red,  $108$ points are colored green, $108$ points are colored blue  and the remaining $108$ points are colored yellow.  Prove that one can choose three points of each color in such a way that the four triangles formed by the chosen points of the same color are congruent.


 



\textbf{Solution}

If you rotate the red points 431 times, they will overlap with blue points $108\times 108$ times, for an average of $\frac{108\times 108}{431}$ per rotation. Note that this average is slightly greater than 27. Therefore at some point 28 red points overlap with blue points. In other words, there exist 28 red and blue points such that the convex 28-gons formed by them are congruent.

 Rotate these 28 red points 431 times. They will overlap with green points $108\times 28$ times, for an average of $\frac{108\times 28}{431}$ per rotation. This average is slightly greater than 7, so at some point 8 of those red points overlap with green points. In other words, there exist 8 red points, 8 blue points, and 8 green points such that the convex octagons formed by them are congruent.

 Rotate these 8 red points 431 times. They will overlap with yellow points $108\times 8$ times, for an average of $\frac{108\cdot 8}{431}$ per rotation. This average is slightly greater than 2, so at some point 3 of those red points will overlap with yellow points. In other words, there exist 3 red points, 3 blue points, 3 green points, and 3 yellow points such that the triangles formed by them are congruent.

 \textbf{Note} Huge bash of Pigeonhole Principle:

 
 \[108\longrightarrow 28\longrightarrow 8\longrightarrow 3\]

 $\left \lceil{\left \lceil{\left \lceil{108\cdot \frac{108}{431}}\right \rceil \cdot \frac{108}{431}}\right \rceil \cdot \frac{108}{431}}\right \rceil =\left \lceil{\left \lceil{28\cdot \frac{108}{431}}\right \rceil \cdot \frac{108}{431}}\right \rceil =\left \lceil{8\cdot \frac{108}{431}}\right \rceil =3$

 


	\item (\textit{USAMO 2012}) Determine which integers $n > 1$ have the property that there exists an infinite sequence $a_1$, $a_2$, $a_3$, $\dots$ of nonzero integers such that the equality

 \[a_k + 2a_{2k} + \dots + na_{nk} = 0\]
holds for every positive integer $k$.


 



\textbf{Solution that involves a non-elementary result}

(Since Bertrand's is well known and provable using elementary techniques, I see nothing wrong with this-tigershark22)

 For $n=2$, $|a_1| = 2 |a_2| = \cdots = 2^m |a_{2^m}|$ implies that for any positive integer $m$, $|a_1| \ge 2^m$, which is impossible.

 We proceed to prove that the infinite sequence exists for all $n\ge 3$.

 First, one notices that if we have $a_{xy} = a_x a_y$ for any integers $x$ and $y$, then it is suffice to define all $a_x$ for $x$ prime, and one only needs to verify the equation (*)

 
 \[a_1+2a_2+\cdots+na_n=0\]

 for the other equations will be automatically true.

 To proceed with the construction, I need the following fact: for any positive integer $m>2$, there exists a prime $p$ such that $\frac{m}{2} <p \le m$.

 To prove this, I am going to use \textbf{Bertrand's Postulate} (\href{https://artofproblemsolving.comhttp://en.wikipedia.org/wiki/Bertrand's\_postulate}{[1]}) without proof. The Theorem states that, for any integer $n>1$, there exists a prime $p$ such that $n<p\le 2n-1$. In other words, for any positive integer $m>2$, if $m=2n$ with $n>1$, then there exists a prime $p$ such that $\frac{m}{2} < p <  m$, and if $m=2n-1$ with $n>1$, then there exists a prime $p$ such that $\frac{m+1}{2} <p\le m$, both of which guarantees that for any integer $m>2$, there exists a prime $p$ such that $\frac{m}{2} <p \le m$.

 


 
Go back to the problem. Suppose $n\ge 3$. Let the largest two primes not larger than $n$ are $P$ and $Q$, and that $n\ge P > Q$. By the fact stated above, one can conclude that $2P > n$, and that $4Q = 2(2Q) \ge 2P > n$. Let's construct $a_n$:

 Let $a_1=1$. There will be three cases: (i) $Q>\frac{n}{2}$,  (ii) $\frac{n}{2} \ge Q > \frac{n}{3}$, and (iii) $\frac{n}{3} \ge Q > \frac{n}{4}$.

 Case (i): $2Q>n$. Let $a_x = 1$ for all prime numbers $x<Q$, and $a_{xy}=a_xa_y$, then (*) becomes:

 
 \[Pa_P + Qa_Q = C_1\]

 Case (ii): $2Q\le n$ but $3Q > n$. In this case, let $a_2=-1$, and $a_x = 1$ for all prime numbers $2<x<Q$, and $a_{xy}=a_xa_y$, then (*) becomes:

 
 \[Pa_P + Qa_Q - Qa_{2Q} = C_2\]

 or

 
 \[Pa_P - Qa_Q = C_2\]

 Case (iii): $3Q\le n$. In this case, let $a_2=3$, $a_3=-2$, and $a_x = 1$ for all prime numbers $3<x<Q$, and $a_{xy}=a_xa_y$, then (*) becomes:

 
 \[Pa_P + Qa_Q + 3Qa_{2Q} - 2Qa_{3Q} = C_3\]

 or
 \[Pa_P + Qa_Q = C_3\]

 In each case, by Bezout's Theorem, there exists non zero integers $a_P$ and $a_Q$ which satisfy the equation. For all other primes $p > P$, just let $a_p=1$ (or any other non-zero integer).

 This construction is correct because, for any $k> 1$, 

 
 \[a_k + 2 a_{2k} + \cdots n a_{nk} =  a_k (1 + 2 a_2 + \cdots n a_n ) = 0\] 

 Since Bertrand's Theorem is not elementary, we still need to wait for a better proof.

 --\href{https://artofproblemsolving.com/wiki/index.php?title=User:Lightest}{Lightest} 21:24, 2 May 2012 (EDT)

 


	\item (\textit{USAMO 2012}) Find all functions $f : \mathbb{Z}^+ \to \mathbb{Z}^+$ (where $\mathbb{Z}^+$ is the set of positive integers) such that $f(n!) = f(n)!$ for all positive integers $n$ and such that $m - n$ divides $f(m) - f(n)$ for all distinct positive integers $m$, $n$.


 



\textbf{Solution 1}

By the first condition we have $f(1)=f(1!)=f(1)!$ and $f(2)=f(2!)=f(2)!$, so $f(1)=1$ or $2$ and similarly for $f(2)$.  By the second condition, we have
 \[n\cdot n!=(n+1)!-n! \mid f(n+1)!-f(n)! \qquad \qquad (1)\]for all positive integers $n$.

 Suppose that for some $n \geq 2$ we have $f(n) = 1$.  We claim that $f(k)=1$ for all $k\ge n$. Indeed, from Equation (1) we have $f(n+1)!\equiv 1 \mod n\cdot n!$, and this is only possible if $f(n+1)=1$; the claim follows by induction.

 We now divide into cases:

 \textbf{Case 1:} $f(1)=f(2)=1$ 

 This gives $f(n)=1$ always from the previous claim, which is a solution. 

 \textbf{Case 2:} $f(1)=2, f(2)=1$ 

 This implies $f(n)=1$ for all $n\ge 2$, but this does not satisfy the initial conditions. Indeed, we would have
 \[3-1 \mid f(3)-f(1)\]and so $2\mid -1$, a contradiction.

 \textbf{Case 3:} $f(1)=1$, $f(2)=2$

 We claim $f(n)=n$ always by induction. The base cases are $n = 1$ and $n = 2$. Fix $k > 1$ and suppose that $f(k)=k$.  By Equation (1) we have that
 \[f(k+1)! \equiv k! \mod k\cdot k! .\]This implies $f(k+1)<2k$ (otherwise $f(k+1)!\equiv 0 \mod k\cdot k!$). Also we have
 \[(k+1)-1  \mid  f(k+1)-f(1)\]so $f(k+1)\equiv 1 \mod k$. This gives the solutions $f(k+1)=1$ and $f(k+1)=k+1$. The first case is obviously impossible, so $f(k + 1) = k + 1$, as desired.  By induction, $f(n) = n$ for all $n$.  This also satisfies the requirements.

 \textbf{Case 4:} $f(1)=f(2)=2$

 We claim $f(n)=2$ by a similar induction. Again if $f(k)=2$, then by (1) we have
 \[f(k+1)\equiv 2 \mod k\cdot k!\]and so $f(k+1)<2k$.  Also note that
 \[k+1-1  \mid  f(k+1)-2\]and
 \[k+1-2  \mid  f(k+1)-2\]so $f(k+1)\equiv 2 \mod k(k-1)$. Then the only possible solution is $f(k+1)=2$. By induction, $f(n) = 2$ for all $n$, and this satisfies all requirements.

 
In summary, there are three solutions: $\boxed{f(n)=1, f(n)=2, f(n)=n}$.

 \textbf{Summary} Three functions: $f(x)=x$ since $x!=x!$, $1!=1$ and $2!=2$, fixed points on the $x!$ function.

 No elegant$\mod$ argument needed; $f(m)-f(n)\textcolor{red}{\textbf{ = }}m-n$ so of course $m-n\mid f(m)-f(n)$!

 



\textbf{Solution 2}

Let $f(n) = c$ and notice $c! = c \implies (c-1)! = 1 \implies c = 1,2$, hence the polynomials $f(x) = 1,2$ work.

 Now, we find all other functions which are nonconstant. Assume $f(x) = a_dx^d + \dots + a_0$. Then, $f(n!) = a_d (n!)^d + \dots + a_0$. Also, $(f(n))! = (f(n))(f(n)-1)(\dots)(1)$. Hence the equality
 \[f(n!) = a_d (n!)^d + \dots + a_0 =  (f(n))(f(n)-1)(\dots)(1)\]holds. The leading coefficients satisfy $a_d = a_d^n \implies a_d = 1$.

 However, the degree of the left side is $d$ and the degree of the right is $d!$ so we have $d = d!$, and that only works if $1 = (d-1)!$, so $d = 1,2$. If $d = 1$, it is easy to see the polynomial $f(x)=x$ works. If $d = 2$, then $(n!)^2 + b(n!)+ c = (n^2+bn+c)(n^2+bn+c-1)(\dots)(1)$. Checking degree however we see that $2 \neq m$ where $m > 2$ is the degree of the RHS. Hence, a contradiction.

 The only three polynomials are from the family 

 
 \[\boxed{f(x) = 1, f(x) = 2, f(x) = x} \quad \forall \ x\].

 It easily follows that $m-n \mid f(m)-f(n) \qquad \square$.

 \textbf{Remark:}

 One can assume that $f(n)$ is a polynomial here rather than an arbitrary function as 

 
                \begin{itemize}
                    \item the statement $f(n!) = (f(n))!$ is nested
                \end{itemize}


                
                \begin{itemize}
                    \item $f(x)$ has the property for all $m,n$, $m-n \mid f(m) - f(n)$.
                \end{itemize}


                The second condition is true for all polynomials, and the \textbf{second uniqueness lemma}, makes the first condition true. If $f(x)$ was not a polynomial, then it follows that both properties cannot hold simultaneously, particularly $m-n \mid f(m) - f(n)$ as I can construct a non-continuous function that satisfies rule 1 but breaks rule 2 or vice versa.

 


	\item (\textit{USAMO 2012}) Let $P$ be a point in the plane of triangle $ABC$, and $\gamma$ a line passing through $P$.  Let $A'$, $B'$, $C'$ be the points where the reflections of lines $PA$, $PB$, $PC$ with respect to $\gamma$ intersect lines $BC$, $AC$, $AB$, respectively.  Prove that $A'$, $B'$, $C'$ are collinear.


 



\textbf{Solution}

By the \href{https://artofproblemsolving.com/wiki/index.php?title=Law\_of\_Sines}{sine law} on triangle $AB'P$,
 \[\frac{AB'}{\sin \angle APB'} = \frac{AP}{\sin \angle AB'P},\]so
 \[AB' = AP \cdot \frac{\sin \angle APB'}{\sin \angle AB'P}.\]

 
                    \begin{figure}[H]
                        \centering
                        \includegraphics[width=0.5\textwidth]{Solution/USAMO_I_2012_5_Solution_Figure_1.png}
                    \end{figure}
                    

 Similarly,
 \begin{align*} B'C &= CP \cdot \frac{\sin \angle CPB'}{\sin \angle CB'P}, \\ CA' &= CP \cdot \frac{\sin \angle CPA'}{\sin \angle CA'P}, \\ A'B &= BP \cdot \frac{\sin \angle BPA'}{\sin \angle BA'P}, \\ BC' &= BP \cdot \frac{\sin \angle BPC'}{\sin \angle BC'P}, \\ C'A &= AP \cdot \frac{\sin \angle APC'}{\sin \angle AC'P}. \end{align*}
Hence,
 \begin{align*} &\frac{AB'}{B'C} \cdot \frac{CA'}{A'B} \cdot \frac{BC'}{C'A} \\ &= \frac{\sin \angle APB'}{\sin \angle AB'P} \cdot \frac{\sin \angle CB'P}{\sin \angle CPB'} \cdot \frac{\sin \angle CPA'}{\sin \angle CA'P} \cdot \frac{\sin \angle BA'P}{\sin \angle BPA'} \cdot \frac{\sin \angle BPC'}{\sin \angle BC'P} \cdot \frac{\sin \angle AC'P}{\sin \angle APC'}. \end{align*}


 Since angles $\angle AB'P$ and $\angle CB'P$ are supplementary or equal, depending on the position of $B'$ on $AC$,
 \[\sin \angle AB'P = \sin \angle CB'P.\]Similarly,
 \begin{align*} \sin \angle CA'P &= \sin \angle BA'P, \\ \sin \angle BC'P &= \sin \angle AC'P. \end{align*}


 By the reflective property, $\angle APB'$ and $\angle BPA'$ are supplementary or equal, so
 \[\sin \angle APB' = \sin \angle BPA'.\]Similarly,
 \begin{align*} \sin \angle CPA' &= \sin \angle APC', \\ \sin \angle BPC' &= \sin \angle CPB'. \end{align*}
Therefore,
 \[\frac{AB'}{B'C} \cdot \frac{CA'}{A'B} \cdot \frac{BC'}{C'A} = 1,\]so by \href{https://artofproblemsolving.com/wiki/index.php?title=Menelaus\%27\_Theorem}{Menelaus's theorem}, $A'$, $B'$, and $C'$ are collinear.

 



\textbf{Solution 2, Barycentric (Modified by Evan Chen)}

We will perform barycentric coordinates on the triangle $PCC'$, with $P=(1,0,0)$, $C'=(0,1,0)$, and $C=(0,0,1)$. Set $a = CC'$, $b = CP$, $c = C'P$ as usual. Since $A$, $B$, $C'$ are collinear, we will define $A = (p : k : q)$ and $B = (p : \ell : q)$.

 Claim: Line $\gamma$ is the angle bisector of $\angle APA'$, $\angle BPB'$, and $\angle CPC'$. This is proved by observing that since $A'P$ is the reflection of $AP$ across $\gamma$, etc.

 Thus $B'$ is the intersection of the isogonal of $B$ with respect to $\angle P$with the line $CA$; that is,
 \[B' = \left( \frac pk \frac{b^2}{\ell}: \frac{b^2}{\ell} : \frac{c^2}{q} \right).\]Analogously, $A'$ is the intersection of the isogonal of $A$ with respect to $\angle P$with the line $CB$; that is,
 \[A' = \left( \frac{p}{\ell} \frac{b^2}{k} : \frac{b^2}{k} : \frac{c^2}{q} \right).\]The ratio of the first to third coordinate in these two pointsis both $b^2pq : c^2k\ell$, so it follows $A'$, $B'$, and $C'$ are collinear.

 



\textbf{Solution 3, Cartesian}

We will use Cartesian coordinates. Without loss of generality, let $P = (0,0)$ and let $\gamma$ be the line $x = 0$. Let $A = (x_1,y_1)$, $B = (x_2,y_2)$, and $C = (x_3,y_3)$. Then $A'$ is the intersection of the lines\begin{align*}y &= -\frac{y_1}{x_1}x \\y - y_2 &= \frac{y_2 - y_3}{x_2 - x_3}(x - x_2).\end{align*}
Solving the system of equations, we see it is
 \[\left(\frac{x_1(x_2y_3 - y_2x_3)}{x_1y_2 + y_1x_2 - x_1y_3 - x_3y_1},\frac{y_1(y_2x_3 - x_2y_3)}{x_1y_2 + y_1x_2 - x_1y_3 - x_3y_1}\right).\]To check collinearity, we need the determinant of the matrix
 \[\begin{bmatrix} x_1(x_2y_3 - y_2x_3) & y_1(y_2x_3 - x_2y_3) & x_1y_2 + y_1x_2 - x_1y_3 - x_3y_1 \\ x_2(x_3y_1 - y_3x_1) & y_2(y_3x_1 - x_3y_1) & x_2y_3 + y_2x_3 - x_2y_1 - x_1y_2 \\ x_3(x_1y_2 - y_1x_2) & y_3(y_1x_2 - x_1y_2) & x_3y_1 + y_3x_1 - x_3y_2 - x_2y_3 \end{bmatrix}\]to be zero, which is true because the columns sum to zero.

 


	\item (\textit{USAMO 2012}) For integer $n \ge 2$, let $x_1$, $x_2$, $\dots$, $x_n$ be real numbers satisfying

 \[x_1 + x_2 + \dots + x_n = 0, \quad \text{and} \quad x_1^2 + x_2^2 + \dots + x_n^2 = 1.\]
For each subset $A \subseteq \{1, 2, \dots, n\}$, define

 \[S_A = \sum_{i \in A} x_i.\]
(If $A$ is the empty set, then $S_A = 0$.)


 Prove that for any positive number $\lambda$, the number of sets $A$ satisfying $S_A \ge \lambda$ is at most $2^{n - 3}/\lambda^2$.  For what choices of $x_1$, $x_2$, $\dots$, $x_n$, $\lambda$ does equality hold?


 



\textbf{Solution 1}

For convenience, let $N=\{1,2,\dots,n\}$.

 Note that $2\sum_{1\leq i<j\leq n} x_ix_j=\left(\sum_{i=1}^{n}x_i\right)^2-\left(\sum_{j=1}^{n} x_i^2\right)=-1$, so the sum of the $x_i$ taken two at a time is $-1/2$. Now consider the following sum:

 
 \[\sum_{A\subseteq N}S_A^2=2^{n-1}\left(\sum_{j=1}^{n} x_i^2\right)+2^{n-1}\left(\sum_{1\leq i<j\leq n} x_ix_j\right)=2^{n-2}.\]

 Since $S_A^2\geq 0$, it follows that at most $2^{n-2}/\lambda^2$ sets $A\subseteq N$ have $|S_A|\geq \lambda$.

 Now note that $S_A+S_{N/A}=0$. It follows that at most half of the $S_A$ such that $|S_A|\geq\lambda$ are positive. This shows that at most $2^{n-3}/\lambda^2$ sets $A\subseteq N$ satisfy $S_A\geq \lambda$.

 Note that if equality holds, every subset $A$ of $N$ has $S_A\in\{-\lambda,0,\lambda\}$. It immediately follows that $(x_1,x_2,\ldots , x_n)$ is a permutation of $(\lambda,-\lambda,0,0,\ldots , 0)$. Since we know that $\sum_{i=1}^{n} x_i^2=1$, we have that $\lambda=1/\sqrt{2}$.

 



\textbf{Solution 2}

Let $x_i=x_1,x_2,...,x_n$ It is evident that $x_i = \frac{(-1)^i}{\sqrt{n}}$ for evens because of the second equation and $x_i=\frac{(-1)^i}{\sqrt{n-1}}$ for odds(one term will be 0 to maintain the first condition). We may then try and get an expression for the maximum number of sets that satisfy this which occur when $\lambda = \frac{1}{\sqrt{n}}$:since it will be 
 \[x_1 + x_2 + ... + x_n\]for any choice of A we pick, it will have to be greater than $\frac{1}{\sqrt{n}}$ which means we can either pick 0 negative $x_m$ or $j-1$ negatives for j positive terms. Since we also have that there are $\frac{n}{2}$ positive and negative terms for evens. Which then gives us:
 \[\sum_{k=1}^{\frac{n}{2}} \binom{\frac{n}{2}}{k}\sum_{i=0}^{k-1} \binom{\frac{n}{2}}{i}\]and 
 \[\sum_{k=1}^{\frac{n}{2}}\binom{\frac{n}{2}}{k}\left(2^{\frac{n}{2}} - \sum_{i=k}^{\frac{n}{2}} \binom{\frac{n}{2}}{i}\right)\]For odd values, let it be the same as the last even valued sequence where n is even(i.e. the same as the sequence before it but with an extra 0 in one of the spots). Then, the following is apparent: 
 \[\sum_{k=1}^{\lfloor\frac{n}{2}\rfloor}\binom{\lfloor \frac{n}{2}\rfloor}{k}\left(2^{\lfloor\frac{n}{2}\rfloor} - \sum_{i=k}^{\lfloor\frac{n}{2}\rfloor}\binom{\lfloor \frac{n}{2}\rfloor}{i}\right)  \le n2^{n-3}\]Thus, we may say that this holds to be true for all $n \ge 2$ since $n2^{n-3}$ grows faster than the sum. Note that equality holds when $S_A \in \{\lambda,0,-\lambda\}$ for all i which occurs when $x_i= \frac{1}{\sqrt{2}},-\frac{1}{\sqrt{2}},0,....$ and $\lambda = \frac{1}{\sqrt{2}}$ since $S_A=\frac{1}{\sqrt{2}}$ is the only choice for $S_A \ge \lambda$$\blacksquare$

 \textbf{Note:} The proof to the last inequality is as follows:First we may rewrite this as being:
 \[\sum_{k=1}^{\frac{n}{2}} \binom{\frac{n}{2}}{k} \left( 2^{\frac{n}{2}} - \sum_{i=k}^{\frac{n}{2}} \binom{\frac{n}{2}}{i}\right)  \le n2^{n-3}\]Thus,
 \[2^n - 2^{\frac{n}{2}} - \sum_{k=1}^{\frac{n}{2}} \binom{\frac{n}{2}}{k} \sum_{i=k}^{\frac{n}{2}} \binom{\frac{n}{2}}{i} \le n2^{n-3}\](the second equation is because the sum of the binomial coefficient is $2^{\frac{n}{2}} - 1$)
 \[2^n - 2^{\frac{n}{2}} \le n2^{n-3} + \sum_{k=1}^{\frac{n}{2}} \binom{\frac{n}{2}}{k} \sum_{i=k}^{\frac{n}{2}} \binom{\frac{n}{2}}{i}\]Since $2^n - 2^{\frac{n}{2}} \le n2^{n-3}$ for all $n \ge 8$ and $\sum_{k=1}^{\frac{n}{2}} \binom{\frac{n}{2}}{k} \sum_{i=k}^{\frac{n}{2}} \binom{\frac{n}{2}}{i} \ge 0$ for all $k$ and $i$, it is apparent that:
 \[\sum_{k=1}^{\frac{n}{2}} \binom{\frac{n}{2}}{k} \sum_{i=0}^{k-1} \binom{\frac{n}{2}}{i} \le n2^{n-3}\]must be true for all $n \ge 8$(because if we rewrite this we get $2^{\frac{n}{2}}\left(8-n\right) \le 8$.) For all $2 \le n < 8$ however, we may use some logic to first layout a plan. Since for $n=6,n=4,$ and $n=2$, $2^{n} - 2^{\frac{n}{2}} = n2^{n-3} + 1$, we may say that whole sum will be less than $n2^{n-3}$ because $\sum_{k=1}^{\frac{n}{2}} \binom{\frac{n}{2}}{k} \sum_{i=k}^{\frac{n}{2}} \binom{\frac{n}{2}}{i} \ge 1$ for all $n \ge 2$ Plugging this inequality back in gives us:
 \[2^{n} - 2^{\frac{n}{2}} - \sum_{k=1}^{\frac{n}{2}} \binom{\frac{n}{2}}{k} \sum_{i=k}^{\frac{n}{2}} \binom{\frac{n}{2}}{i} \le 2^{n} - 2^{\frac{n}{2}} - 1 = n2^{n-3}\] because of the fact that $- \sum_{k=1}^{\frac{n}{2}} \binom{\frac{n}{2}}{k} \sum_{i=k}^{\frac{n}{2}} \binom{\frac{n}{2}}{i} \le -1$

 


\end{enumerate}